\section*{Erd\H{o}s Problem 80: books in graphs whose edges all lie in triangles}
\subsection*{1. Statement and scope}
For fixed $0<c<1/2$, let $f_c(n)$ be the least possible maximum number of triangles through an edge in an $n$-vertex graph with at least $cn^2$ edges and every edge in a triangle. The restriction on $c$ ensures that admissible graphs exist for large $n$. We prove elementary density bounds and deduce $f_c(n)\to\infty$ from the triangle removal lemma, stating that input precisely. The finer rate for $c<1/4$, including the proposed logarithmic lower bound, remains unresolved here.

\subsection*{2. Two counts and an elementary linear bound}
Let $m=e(G)$, $T$ be its number of triangles, and $B$ its maximum book size. For an edge $uv$, its triangle count is $|N(u)\cap N(v)|$. Counting each triangle at its three edges gives
\[
 m\le3T\le mB. \tag{1}
\]
The first inequality uses the hypothesis that every edge is covered. Also
\[
 \sum_{uv\in E(G)}(d(u)+d(v))=\sum_v d(v)^2\ge4m^2/n.
\]
Some edge has endpoint-degree sum at least $4m/n$. Since $N(u)\cup N(v)$ has at most $n$ vertices, that edge has at least $4m/n-n$ common neighbours. Thus
\[
 B\ge (4c-1)n\quad\text{if }m\ge cn^2. \tag{2}
\]
This is informative for $c>1/4$; it does not give a positive lower bound in the main sparse-density regime.

\subsection*{3. Qualitative divergence via triangle removal}
Use the standard triangle removal lemma as an external input: for every $\eta>0$ there is $\delta>0$ such that an $n$-vertex graph with fewer than $\delta n^3$ triangles can be made triangle-free by deleting at most $\eta n^2$ edges.
If $T\ge\delta n^3$, then (1) and $m\le n^2/2$ give $B\ge6\delta n$. Otherwise, let $D$ be a set of at most $\eta n^2$ edges whose deletion kills every triangle. Each original triangle contains an edge of $D$, and each edge of $D$ lies in at most $B$ original triangles. Therefore $T\le B|D|$, and (1) gives
\[
 cn^2\le m\le3T\le3B\eta n^2.
\]
In all cases,
\[
 B\ge\min\{6\delta(\eta)n,\ c/(3\eta)\}. \tag{3}
\]
Given any $M>0$, choose $\eta=c/(6(M+1))$. Once $n> (M+1)/(6\delta(\eta))$, both entries on the right of (3) exceed $M$. This proves $f_c(n)\to\infty$ with the correct order of quantifiers: $\eta$ and its removal constant are fixed before $n$ grows. The proof supplies no good explicit rate without an appropriate quantitative estimate for $\delta(\eta)$.

\subsection*{4. A check on the density bound}
For $n$ divisible by three, the complete tripartite graph with three equal parts has $m=n^2/3$, and every edge has exactly $n/3$ common neighbours, all in the third part. Hence $f_c(n)\le n/3$ for $c\le1/3$ on this subsequence. At $c=1/3$, (2) matches this upper bound exactly: $f_{1/3}(n)=n/3$ when $3\mid n$.
This special exact value does not estimate the sublinear regime $c<1/4$. The problem page records the Fox--Loh subpolynomial upper construction there, so a fixed positive power lower bound in that regime is already false; that construction is background, not proved in this text.

\subsection*{5. Assessment and attribution}
The counting argument, exact tripartite case, and reduction of divergence to the stated removal lemma are complete. The logarithmic-rate question is unresolved in this attempt. This replaces earlier GPT-5.2 write-ups with an explicit removal argument and carefully scoped conclusions, without claiming a new bound or unperformed exhaustive search.
Sources: \url{https://www.erdosproblems.com/80}; J. Fox, \emph{A new proof of the graph removal lemma}, \url{https://arxiv.org/abs/1006.1300}.

\textbf{Completion rate estimate: 25\%.}
